% SPDX-License-Identifier: LPPL-1.3c
% Copyright (C) 2026 Christophe Jorssen
% Author and Current Maintainer: Christophe Jorssen <christophe.jorssen@gmail.com>
% LPPL maintenance status: maintained
% This file is part of luafemm. See LICENSE and LICENSES.md for its terms.
\section{Three-dimensional fields of prescribed currents}
\label{sec:biot}

A loop drawn in a plane still produces a three-dimensional magnetic field.
The |femm/biot| provider evaluates that field from prescribed steady currents
in a homogeneous linear medium. It needs no finite-element mesh. Filaments,
round volume-current wires, circular and arbitrary planar loops, straight
coils and toroidal windings share one model and may be superposed.

Use |femm/problem| for induced magnetization in iron, permanent-magnet material
laws or resolved eddy currents. The prescribed-current provider does not solve
for a redistribution of current, skin effect, conductivity, magnetic interfaces,
forces or inductance. A camera projection is a presentation of a 3D field, not
a planar or axisymmetric finite-element formulation.

\subsection{A loop and its axial field}
Create the model at picture initialization, declare all sources, and then capture
profiles or draw fields. A field request freezes the declarations. The following
example compares the computed field with
\[
 B_z(z)=\frac{\mu_0 I R^2}{2(R^2+z^2)^{3/2}}.
\]
The model uses centimetres, so $R=1$ means $0.01\,\mathrm m$. The current is in
amperes and the returned field in teslas. The same source supplies the diagram
and the profile. The camera changes neither of them physically.
\begin{codeexample}[width=0pt]
\begingroup
\tikzset{source view/.style={femm/biot,
 femm/biot view={azimuth=-30,elevation=25,step=.12,
 maximum length=20,xmin=-2.5,xmax=2.5,ymin=-2.5,ymax=2.5,
 zmin=-2.5,zmax=2.5}}}
\tikzpicture[source view]
 \femmbiotsource[kind=loop,current=2,radius=1]
 \femmbiotprofile[name=biot-axis,from=0/0/-2,to=0/0/2,samples=61]
 \pic {femm biot field};
 \pic {femm biot geometry};
\endtikzpicture
\par
\tikzpicture
 \axis[width=8cm,height=5cm,no markers,grid=major,
 xlabel={$z$ (cm)},ylabel={$B_z$ (T)},legend style={at={(.5,-.32)},anchor=north,legend columns=2}]
  \femmbiotplot[profile=biot-axis,component=Bz,abscissa=z,
   plot options={blue!70!black,thick}]
  \addlegendentry{Biot--Savart}
  \addplot[orange!80!black,dashed,domain=-2:2,samples=101]
   {4*pi*1e-5/(1+x*x)^1.5};
  \addlegendentry{On-axis formula}
 \endaxis
\endtikzpicture
\par
\endgroup
\end{codeexample}

\subsection{A wire with a finite radius}
An infinite round wire of radius $a$ carrying a uniform current density has
$B_\theta=\mu I r/(2\pi a^2)$ inside and $B_\theta=\mu I/(2\pi r)$ outside.
The nonzero |wire radius| selects the volume-current model described below.
The diagram is a transverse section, and the signed profile crosses the axis.
A missing filament sample at the axis produces a gap, never a zero-field value.
The filled disk is ordinary TikZ annotation of the physical radius; the generated
3D source picture itself draws centre lines.
\begin{codeexample}[width=0pt]
\begingroup
\tikzpicture[femm/biot,
 femm/biot view={azimuth=0,elevation=90,xmin=-1.5,xmax=1.5,
 ymin=-1.5,ymax=1.5,zmin=-1,zmax=1,
 seeds={.2/0/0;.6/0/0;1.1/0/0},step=.08}]
 \femmbiotsource[kind=wire,extent=infinite,current=3,wire radius=.4]
 \femmbiotprofile[name=thick-wire,from=-1.4/0/0,to=1.4/0/0,samples=141]
 \filldraw[fill=orange!15,draw=orange!75!black] (0,0) circle[radius=.4];
 \pic {femm biot field};
 \node at (0,0) {$\odot$};
\endtikzpicture
\tikzpicture[femm/biot]
 \femmbiotsource[kind=wire,extent=infinite,current=3]
 \femmbiotprofile[name=thin-wire,from=-1.4/0/0,to=1.4/0/0,samples=141]
\endtikzpicture
\par
\tikzpicture
 \axis[width=8cm,height=5cm,no markers,grid=major,
 xlabel={$x$ (cm)},ylabel={$B_y$ (T)},ymin=-.0003,ymax=.0003,
 legend pos=north west]
  \femmbiotplot[profile=thick-wire,component=By,abscissa=x,
   plot options={blue!70!black,thick}]
  \addlegendentry{Round wire}
  \femmbiotplot[profile=thin-wire,component=By,abscissa=x,
   plot options={orange!80!black,dashed}]
  \addlegendentry{Filament}
 \endaxis
\endtikzpicture
\par
\endgroup
\end{codeexample}

\subsection{Polygonal and Bezier loops}
A |femm/current path| postaction captures a connected evaluated PGF path.
Use ordinary coordinates, node anchors, straight segments, arcs, rounded corners
or cubic Bezier controls. The source is defined in a chosen physical plane,
then rotated and translated in 3D. The original path need not itself be drawn:
|femm biot geometry| supplies its camera projection.

PGF paths are two-dimensional after evaluation. A TikZ coordinate with three
components has already been projected by PGF and cannot be inverted into a
unique 3D curve. Use |plane|, |offset| and the source pose for planar paths;
use explicit |from|/|to| triples or a generated helix for spatial sources.
Circular |kind=loop| sources use an analytic circle; a TikZ |circle| captured
as a path retains PGF's Bezier approximation and the additional chord tolerance.
\begin{codeexample}[width=0pt]
\begingroup
\tikzpicture[femm/biot,femm/biot view={seeds={0/0/0;.4/0/0},
 step=.15,maximum length=15}]
 \path[femm/current path={current=2,curve tolerance=.001}]
  (-1,-.7)--(1,-.7)..controls(1.6,.9) and(-1.6,.9)..cycle;
 \pic {femm biot field};
 \pic {femm biot geometry};
\endtikzpicture
\qquad
\tikzpicture[femm/biot,femm/biot view={seeds={0/0/0;.4/0/0},
 step=.15,maximum length=15}]
 \path[femm/current path={current=2,rotate y=25}]
  (-1,-1)--(1,-1)--(1,1)--(-1,1)--cycle;
 \pic {femm biot field};
 \pic {femm biot geometry};
\endtikzpicture
\par
\endgroup
\end{codeexample}

\subsection{Independent loops and a continuous helical winding}
For |winding=turns|, the straight coil contains independent circular loops,
equally spaced from $-L/2$ to $L/2$. A single turn is placed at zero. They may be
interpreted as a series winding whose external leads have been omitted from this
idealized source model; each included loop is individually closed.

For |winding=helix|, the path makes $N$ full revolutions while advancing by $L$.
An axial return adds two radial leads and a straight return along the axis.
Their field is included. This geometry is an explicit example, not an automatic
model of an experimental feed cable. To combine an open contribution with a
separately described feed, request |closure=contribution| and supply the rest
of the circuit. Never interpret isolated finite endpoints as a complete steady
current distribution.
\begin{codeexample}[width=0pt]
\begingroup
\pgfkeys{/femm/biot source/coil sources/.style={current=2,radius=1,length=3,
 turns=6,wire radius=.04}}
\tikzset{coil view/.style={femm/biot,
 femm/biot view={azimuth=-20,elevation=20,
 seeds={0/.35/0;0/-.35/0},step=.15,maximum length=16,
 geometry samples=32}}}
\tikzpicture[coil view]
 \femmbiotsource[coil sources,kind=solenoid,winding=turns]
 \pic {femm biot field};\pic {femm biot geometry};
 \node[below=3mm] at (current bounding box.south) {Independent loops};
\endtikzpicture
\qquad
\tikzpicture[coil view]
 \femmbiotsource[coil sources,kind=solenoid,winding=helix,closure=axis]
 \pic {femm biot field};\pic {femm biot geometry};
 \node[below=3mm] at (current bounding box.south) {Helix and axial return};
\endtikzpicture
\par
\endgroup
\end{codeexample}

\subsection{Toroidal windings}
Independent poloidal loops lie in meridional planes at uniformly spaced azimuths.
The toroidal helix makes one major-circle revolution and $N$ poloidal revolutions;
it is closed for the required integer turn count. Its additional toroidal current
component changes the field, so these two windings need not have identical
leakage or field lines. A finite number of turns is not the perfectly confined
continuous-sheet toroid. Overlapping volume-current supports simply superpose;
this provider does not check insulation or winding manufacture.
\begin{codeexample}[width=0pt]
\begingroup
\pgfkeys{/femm/biot source/toroidal winding/.style={current=2,
 radius=.55,major radius=1.7,turns=12,wire radius=.03}}
\tikzset{toroid view/.style={femm/biot,
 femm/biot view={azimuth=-35,elevation=45,
 seeds={1.7/0/0;1.7/0/.25},step=.15,
 maximum length=12,geometry samples=32}}}
\tikzpicture[toroid view]
 \femmbiotsource[toroidal winding,kind=toroid,winding=turns]
 \pic {femm biot geometry};\pic {femm biot field};
 \node[below=3mm] at (current bounding box.south) {Poloidal loops};
\endtikzpicture
\qquad
\tikzpicture[toroid view]
 \femmbiotsource[toroidal winding,kind=toroid,winding=helix]
 \pic {femm biot geometry};\pic {femm biot field};
 \node[below=3mm] at (current bounding box.south) {Toroidal helix};
\endtikzpicture
\par
\endgroup
\end{codeexample}

\subsection{Commands, scope and operation order}
\begin{key}{/tikz/femm/biot=\meta{options} (default empty)}
Initialize a prescribed-current picture using |/femm/biot|. Use a separate picture
from the FE and electric providers. The default drawing basis is one centimetre
per model coordinate, independently of its physical unit. A root TikZ transform
changes display only. All numerical options below accept finite numeric literals,
including scientific notation; ordinary TikZ coordinate expressions are evaluated
by PGF before a path is captured.
\end{key}
\begin{command}{\femmbiotsource\oarg{options}}
Declare one source using |/femm/biot source|. Options are grouped and may be
composed with ordinary PGF styles in that family. Supply the current explicitly.
Position and orientation use the source pose keys, independently of the camera
and local drawing scopes. At most 2000 source declarations and 100000 initial curve panels are accepted.
\end{command}
\begin{key}{/tikz/femm/biot source=\meta{options}}
Invoke the same declaration through a TikZ key, after model initialization.
\end{key}
\begin{key}{/tikz/femm/current path=\meta{options}}
Capture one connected physical path in a postaction. Options belong to
|/femm/biot source|; |kind=wire| is selected by the capture. Use |--cycle| or an
explicitly returning path for closed circuits. Local TikZ path transforms affect
its evaluated planar coordinates. Do not draw fields or save the model among
this path's own options: its capture has not yet occurred there.
\end{key}
\begin{key}{/tikz/femm/biot view=\meta{options}}
Set grouped display/tracing controls from |/femm/biot view|. Declare all physical
sources before the first picture, profile or value request. Drawing geometry
also freezes the source list. Subsequent views and profiles remain available.
\end{key}
\begin{key}{/tikz/pics/femm biot geometry}
Draw projected source centre lines. An infinite wire is clipped to the physical
view box. Curved geometry uses the independent |geometry samples| control.
This drawing is not a surface rendering of a thick conductor.
\end{key}
\begin{key}{/tikz/pics/femm biot field}
Integrate true 3D field lines in both directions by default, then project their
vertices. Every arrow follows $\mathbf B$, including a backward-only request.
Source exclusions, null fields, the physical box, length and step budgets stop
integration. A near-return to the seed closes a curve at the tracing resolution.
Line counts and seed placement are aesthetic and do not measure flux.
\end{key}
\begin{key}{/tikz/pics/femm biot vectors}
Sample the full 3D field on a selected coordinate plane, normalize in 3D and
project each arrow. A vector normal to the screen can disappear in projection;
it is not a zero magnetic field. No in-plane streamline approximation is used.
\end{key}
\begin{key}{/tikz/femm biot conductor}
Default orange-brown centre line, width |.9pt|. Ordinary TikZ styles control
colour, thickness and annotations independently of physical radius.
\end{key}
\begin{key}{/tikz/femm biot field}
Default dark-blue line, width |.4pt|, with a Stealth arrow at relative position
|.55|. Use |/.append style| to preserve the orientation decoration.
\end{key}
\begin{key}{/tikz/femm biot vector}
Default dark-teal arrow, width |.4pt|, with a |1.4mm| Stealth tip.
\end{key}
\begin{command}{\femmbiotprofile\oarg{options}}
Capture an owned 3D polyline profile using |/femm/biot profile|. Samples are
equispaced in physical path length. Names have document lifetime; reusing one
replaces that profile. Later source or picture changes do not alter its samples.
\end{command}
\begin{command}{\femmbiotplot\oarg{options}}
Inside a PGFPlots axis, plot a stored profile using |/femm/biot plot|. It behaves
like |\addplot+|, retaining axis cycle lists and legends. Invalid filament samples
create gaps. Axes, units, legends and formatting remain ordinary PGFPlots options.
\end{command}
\begin{command}{\femmbiotvalue\oarg{options}}
Print one SI field component using |/femm/biot value| and the current PGF number
format. A point in an excluded filament neighbourhood is an error.
\end{command}

\subsection{Model controls}
\begin{key}{/femm/biot/length unit=m\textbar cm\textbar mm\textbar um (initially cm)}
Physical length unit for every source, profile and tracing coordinate. It does not set the printed axis label or change the camera scale.
\end{key}
\begin{key}{/femm/biot/relative permeability=\meta{value} (initially 1)}
Positive, spatially uniform relative permeability. There are no material interfaces or nonlinear laws in this provider.
\end{key}
\begin{key}{/femm/biot/quadrature tolerance=\meta{value} (initially 1e-5)}
Positive relative acceptance tolerance, smaller than one, for each curve integral. Combined with the absolute tolerance; this is a numerical estimate, not a certified error bound.
\end{key}
\begin{key}{/femm/biot/absolute tolerance=\meta{value} (initially 1e-12)}
Positive absolute field tolerance in teslas, shared between source curves at each sample. It is essential near cancellations and nulls.
\end{key}
\begin{key}{/femm/biot/maximum quadrature depth=\meta{value} (initially 20)}
Subdivision depth per initial parametric panel, an integer from 1 to 30. Exhaustion is an error, never a silently accepted field.
\end{key}
\begin{key}{/femm/biot/maximum evaluations=\meta{value} (initially 200000)}
Total integrand evaluations per field sample, an integer from 24 to 10000000. Exact straight-wire formulae need no quadrature evaluations.
\end{key}
\begin{key}{/femm/biot/exclusion radius=\meta{value} (initially .001)}
Positive distance in model units around a zero-radius filament. It masks singular positions without changing the field law elsewhere. It does not apply to a positive-radius volume current.
\end{key}
\begin{key}{/femm/biot/native file=\meta{value} (initially empty)}
Read a prescribed-current |.femx| or |.ansx| at picture initialization. The file owns physical and numerical settings; explicit overrides are rejected. Camera settings remain independent.
\end{key}
\begin{key}{/femm/biot/name=\meta{value} (initially empty)}
Optional document-wide snapshot name, captured at picture end. Use it in later profile/value requests or |\femmsave[model=...]|. Keep magnetic model names unique.
\end{key}
\subsection{Source controls}
\begin{key}{/femm/biot source/kind=wire\textbar loop\textbar solenoid\textbar toroid (initially loop)}
Registered source geometry. A wire uses a straight segment or captured planar path; a loop is an analytic circle; the two coils use the selected winding construction.
\end{key}
\begin{key}{/femm/biot source/current=\meta{value} (initially empty)}
Required signed current in amperes, per loop or along the continuous wire. A positive current follows the source parameter; a negative value reverses the magnetic field.
\end{key}
\begin{key}{/femm/biot source/wire radius=\meta{value} (initially 0)}
Nonnegative physical radius in model units. Zero selects a filament; a positive value selects the round convolution model below. It must be smaller than the winding radius for generated curved sources.
\end{key}
\begin{key}{/femm/biot source/center=\meta{value} (initially 0/0/0)}
Translation triple |x/y/z| in model units, applied after the source rotations. For paths, it translates the lifted evaluated coordinates.
\end{key}
\begin{key}{/femm/biot source/rotate x=\meta{value} (initially 0)}
Rotation in degrees about the source x axis. Rotations are applied in x, then y, then z order, followed by translation.
\end{key}
\begin{key}{/femm/biot source/rotate y=\meta{value} (initially 0)}
Rotation in degrees about the y axis after the x rotation. It rotates geometry and current together.
\end{key}
\begin{key}{/femm/biot source/rotate z=\meta{value} (initially 0)}
Final rotation in degrees about the z axis, before translation.
\end{key}
\begin{key}{/femm/biot source/closure=closed\textbar contribution\textbar axis (initially closed)}
A finite wire must return to its initial point in closed mode. Contribution mode explicitly permits an incomplete circuit. Axis mode is for a straight helical coil and includes its radial leads and axial return.
\end{key}
\begin{key}{/femm/biot source/winding=turns\textbar helix (initially turns)}
For solenoid and toroid sources: independent planar circles or a continuous spatial helix. A straight helix requires an explicit axis return or contribution declaration.
\end{key}
\begin{key}{/femm/biot source/handedness=right\textbar left (initially right)}
Sign of the azimuthal advance of a straight helix, or poloidal advance of a toroidal helix. Current sign remains independent. Applies to helical windings.
\end{key}
\begin{key}{/femm/biot source/extent=finite\textbar infinite (initially finite)}
For wire sources only. Infinite extent requires exactly two distinct points defining the oriented axis. A captured multi-segment path is finite.
\end{key}
\begin{key}{/femm/biot source/from=\meta{value} (initially 0/0/-1)}
Initial triple |x/y/z| in model units for an explicitly declared wire. Source pose is applied afterwards.
\end{key}
\begin{key}{/femm/biot source/to=\meta{value} (initially 0/0/1)}
Final triple |x/y/z| for a wire. With infinite extent, the two points set only the line and positive-current direction, not endpoints.
\end{key}
\begin{key}{/femm/biot source/radius=\meta{value} (initially 1)}
Positive loop radius, straight-coil radius or toroidal minor winding radius, in model units.
\end{key}
\begin{key}{/femm/biot source/major radius=\meta{value} (initially 2)}
Toroidal centre-line radius in model units; it must exceed the minor winding radius.
\end{key}
\begin{key}{/femm/biot source/length=\meta{value} (initially 3)}
Positive axial span of a straight coil in model units. Planar turns include both endpoints of the span when there is more than one turn.
\end{key}
\begin{key}{/femm/biot source/turns=\meta{value} (initially 10)}
Integer turn count from 1 to 1000 for either coil. A single straight-coil planar turn lies at its centre. A toroidal helix has exactly one major revolution.
\end{key}
\begin{key}{/femm/biot source/phase=\meta{value} (initially 0)}
Geometrical starting angle in degrees. For a straight coil it shifts the initial circular angle; for a toroid it shifts the major azimuth. It is not an AC phasor angle.
\end{key}
\begin{key}{/femm/biot source/plane=xy\textbar xz\textbar yz (initially xy)}
Plane used to lift a captured two-dimensional PGF path. Positive planar orientation follows the ordered coordinates: x/y, x/z or y/z; the x/z normal is therefore negative y.
\end{key}
\begin{key}{/femm/biot source/offset=\meta{value} (initially 0)}
Coordinate perpendicular to the captured path plane, in model units, before the source pose is applied.
\end{key}
\begin{key}{/femm/biot source/curve tolerance=\meta{value} (initially .002)}
Positive chord tolerance in model units for PGF path flattening. Refine it independently of quadrature tolerance. Generated analytic circles/helices do not use this polygonalization.
\end{key}
\subsection{View controls}
\begin{key}{/femm/biot view/azimuth=\meta{value} (initially -45)}
Camera angle in degrees. The screen horizontal basis is $(\cos\alpha,\sin\alpha,0)$. No physical source rotation is implied.
\end{key}
\begin{key}{/femm/biot view/elevation=\meta{value} (initially 25)}
Camera angle in degrees. The screen vertical basis is $(-\sin\alpha\sin\beta,\cos\alpha\sin\beta,\cos\beta)$. With azimuth zero and elevation 90 degrees, the screen is the physical x/y plane.
\end{key}
\begin{key}{/femm/biot view/step=\meta{value} (initially .08)}
Maximum Runge--Kutta arc-length step in model units. Adaptive steps may be smaller. It also sets the near-return resolution for closing a traced loop.
\end{key}
\begin{key}{/femm/biot view/tolerance=\meta{value} (initially .0002)}
Positive local streamline error tolerance in model units, estimated by step doubling. It does not change Biot--Savart quadrature or source geometry.
\end{key}
\begin{key}{/femm/biot view/maximum length=\meta{value} (initially 30)}
Positive traced length in model units, separately for each direction from a seed.
\end{key}
\begin{key}{/femm/biot view/maximum steps=\meta{value} (initially 3000)}
Accepted-step budget per direction, an integer from 1 to 100000. Exhaustion retains the partial trace with a diagnostic.
\end{key}
\begin{key}{/femm/biot view/null field=\meta{value} (initially 1e-14)}
Nonnegative threshold in teslas below which a field direction is undefined for tracing or arrows.
\end{key}
\begin{key}{/femm/biot view/geometry samples=\meta{value} (initially 64)}
Integer from 8 to 1000, per circle or minor helical turn in the source drawing. It affects presentation only.
\end{key}
\begin{key}{/femm/biot view/xmin=\meta{value} (initially -3)}
Physical x coordinate of the lower view-box face, in model units. Bounds must be ordered. They stop traces and clip infinite-wire drawings; they impose no magnetic boundary condition.
\end{key}
\begin{key}{/femm/biot view/xmax=\meta{value} (initially 3)}
Physical x coordinate of the upper view-box face, in model units. Bounds must be ordered. They stop traces and clip infinite-wire drawings; they impose no magnetic boundary condition.
\end{key}
\begin{key}{/femm/biot view/ymin=\meta{value} (initially -3)}
Physical y coordinate of the lower view-box face, in model units. Bounds must be ordered. They stop traces and clip infinite-wire drawings; they impose no magnetic boundary condition.
\end{key}
\begin{key}{/femm/biot view/ymax=\meta{value} (initially 3)}
Physical y coordinate of the upper view-box face, in model units. Bounds must be ordered. They stop traces and clip infinite-wire drawings; they impose no magnetic boundary condition.
\end{key}
\begin{key}{/femm/biot view/zmin=\meta{value} (initially -3)}
Physical z coordinate of the lower view-box face, in model units. Bounds must be ordered. They stop traces and clip infinite-wire drawings; they impose no magnetic boundary condition.
\end{key}
\begin{key}{/femm/biot view/zmax=\meta{value} (initially 3)}
Physical z coordinate of the upper view-box face, in model units. Bounds must be ordered. They stop traces and clip infinite-wire drawings; they impose no magnetic boundary condition.
\end{key}
\begin{key}{/femm/biot view/seeds=\meta{value} (initially auto)}
Use |auto| for aesthetic seeds near the first source, or a semicolon-separated list of triples such as |0/0/0;.5/0/0|. At most 500 seeds are allowed, all inside the tracing box. Explicit seeds are preferable for comparisons.
\end{key}
\begin{key}{/femm/biot view/direction=both\textbar forward\textbar backward (initially both)}
Direction of integration from each seed. Emitted paths always carry arrows in the physical field direction.
\end{key}
\begin{key}{/femm/biot view/plane=xy\textbar xz\textbar yz (initially xz)}
Coordinate plane on which full-vector arrows are sampled. It does not constrain the three-dimensional field-line integrator.
\end{key}
\begin{key}{/femm/biot view/offset=\meta{value} (initially 0)}
Perpendicular coordinate of the vector sampling plane, in model units.
\end{key}
\begin{key}{/femm/biot view/spacing=\meta{value} (initially .5)}
Positive vector-grid spacing in model units. At most 10000 sample positions are accepted.
\end{key}
\begin{key}{/femm/biot view/vector length=\meta{value} (initially .25)}
Positive normalized arrow length in 3D model units before projection. Apparent projected lengths may therefore differ.
\end{key}
\subsection{Profile controls}
\begin{key}{/femm/biot profile/name=\meta{value} (initially empty)}
Required document-wide profile identifier: letters/digits followed by letters, digits, colon, underscore or hyphen.
\end{key}
\begin{key}{/femm/biot profile/model=\meta{value} (initially empty)}
Empty selects the current Biot--Savart picture; otherwise select a named snapshot.
\end{key}
\begin{key}{/femm/biot profile/from=\meta{value} (initially 0/0/-3)}
Initial profile point |x/y/z| in model units. It is independent of the camera and source pose.
\end{key}
\begin{key}{/femm/biot profile/to=\meta{value} (initially 0/0/3)}
Final profile point |x/y/z| in model units.
\end{key}
\begin{key}{/femm/biot profile/points=\meta{value} (initially empty)}
Optional semicolon-separated polyline triples. When supplied, replaces from/to. Successive vertices must be distinct.
\end{key}
\begin{key}{/femm/biot profile/samples=\meta{value} (initially 101)}
Integer from 2 to 10001, including both profile endpoints. Samples are uniform in polyline arc length.
\end{key}
\subsection{Plot controls}
\begin{key}{/femm/biot plot/profile=\meta{value} (initially empty)}
Required previously captured profile name. No field recalculation occurs when plotting it.
\end{key}
\begin{key}{/femm/biot plot/component=B\textbar Bx\textbar By\textbar Bz\textbar H\textbar Hx\textbar Hy\textbar Hz (initially B)}
SI field magnitude or Cartesian component. B quantities are in teslas; H quantities are in amperes per metre.
\end{key}
\begin{key}{/femm/biot plot/abscissa=s\textbar x\textbar y\textbar z (initially s)}
Arc length s from the profile start, or a Cartesian coordinate. Values use the model length unit, not necessarily metres.
\end{key}
\begin{key}{/femm/biot plot/abscissa scale=\meta{value} (initially 1)}
Positive multiplier applied to the displayed abscissa, for example 10 to convert a centimetre profile to millimetres. Set the axis label yourself.
\end{key}
\begin{key}{/femm/biot plot/plot options=\meta{value} (initially empty)}
Unexpanded native PGFPlots options for this plot, including reusable styles. They supplement the normal cycle-list choice.
\end{key}
\subsection{Value controls}
\begin{key}{/femm/biot value/model=\meta{value} (initially empty)}
Empty selects the current Biot--Savart picture; otherwise select a named snapshot.
\end{key}
\begin{key}{/femm/biot value/at=\meta{value} (initially 0/0/0)}
Sample triple |x/y/z| in model units. The result is evaluated before PGF number formatting.
\end{key}
\begin{key}{/femm/biot value/component=B\textbar Bx\textbar By\textbar Bz\textbar H\textbar Hx\textbar Hy\textbar Hz (initially B)}
SI field magnitude or Cartesian component. B quantities are in teslas; H quantities are in amperes per metre.
\end{key}

\subsection{Biot--Savart integration and finite-radius currents}
For a filament curve $\mathbf c(t)$, the evaluator computes
\[
 \mathbf B(\mathbf x)=\frac{\mu_0\mu_r I}{4\pi}
 \int\frac{\mathbf c'(t)\mathbin\times(\mathbf x-\mathbf c(t))}
 {|\mathbf x-\mathbf c(t)|^3}\,dt.
\]
Straight segments and infinite wires use exact formulas where applicable.
Analytic curves use nested 4/8-point Gauss--Legendre rules, with initial panels
resolving each turn. The largest estimated contribution is refined first.
Absolute and relative tolerances are checked on each curve integral; summed
error estimates do not certify relative accuracy after cancellation of large
source fields. Budgets fail explicitly. Refine source chord tolerance,
quadrature tolerance, tracing tolerance and profile sampling independently.

For a nonzero radius $a$, the prescribed current is the convolution of the
filament current with the normalized radial kernel
\[
 K_a(r)=\begin{cases}
 [\pi^2 a^2\sqrt{a^2-r^2}]^{-1},&r<a,\\
 0,&r\geq a.
 \end{cases}
\]
For a straight infinite wire, integrating this kernel along the wire gives
exactly $J=I/(\pi a^2)$ throughout its circular cross-section. Convolution
preserves zero divergence for a closed current path. In a bend, it defines
its own curvature-dependent volume-current distribution; it is \emph{not}
an exact solution for uniform current in every normal disk of a curved metal
wire. Finite segment ends are rounded by the same kernel. No conductivity,
charge redistribution, skin or proximity law is solved.

Radial symmetry reduces the volume integration to the same curve integral,
with the factor $1/d^3$ replaced by $F(d/a)/d^3$, where
\[
 F(q)=\frac2\pi\bigl(\arcsin q-q\sqrt{1-q^2}\bigr)\quad(0\leq q<1),
 \qquad F(q)=1\quad(q\geq1).
\]
The removable limit is used at zero. Outside the radius-a neighbourhood of all
filaments, the field equals the filament field for this particular model.
Thus increasing the radius is a physical change to the interior source model,
not a universal approximation to the field of an arbitrary thick curved wire.

Sources and projection data stay separate. The orthographic renderer does not
remove hidden lines, draw occluding conductor surfaces or infer a flux density
from the number of field lines. A projected crossing need not be a physical
intersection in three dimensions.

\subsection{Native storage and scope}
Use |\femmsave| after source capture, with |kind=problem| and a |.femx| filename,
or |kind=answer| and a |.ansx| filename. The registered
|magnetic-biot-1| adapter owns source geometry, currents, units, finite-radius
law and quadrature controls. Answers store the registered evaluation law rather
than fabricated finite-element coefficients. Reading performs no field
integration; a later view or profile evaluates the retained sources.

No FEMM |.fem|/|.ans| representation is provided for these 3D sources. There is
no general magnetic 3D finite-element solve, imported 3D dataset adapter or
material response in this provider. A positive homogeneous permeability
multiplies the Biot--Savart field; it does not model an inserted ferromagnetic core.

\begin{codeexample}[width=0pt]
\begin{tikzpicture}[femm/biot]
 \femmbiotsource[kind=loop,current=1]
 \femmsave[kind=answer,file=manual-biot.ansx]
 \pic {femm biot geometry};
\end{tikzpicture}
\qquad
\begin{tikzpicture}[femm/biot={native file=manual-biot.ansx}]
 \pic[femm/biot view={seeds={0/0/0},maximum length=3}]
  {femm biot field};
 \pic {femm biot geometry};
\end{tikzpicture}
\end{codeexample}
